% ============================================
% Supplementary Note 5: Alternative Models and Controls for Community-Level Selection
% ============================================
\subsection*{\rev{Supplementary Note 5: Alternative models and controls for community-level selection}}
\label{sec:suppl:note5}

\subsubsection*{\rev{Biomass loading}}
\rev{Several alternative mechanisms could generate parental-affiliation-correlated persistence during coalescence~\citep{Mansour2018,Rillig2015}, including passive biomass loading, initial-richness effects, and environmental modification by dominant taxa. Because coalescence was initiated by equal-volume pooling without OD normalization, parental-community biomass differences at mixing could bias the coalesced community toward the higher-biomass parent. To test this hypothesis, we asked whether parental-community OD$_{600}$ correlates with parental dominance direction or with pairwise selection correlation. If Dominance were driven primarily by unequal absolute biomass at mixing, the higher-OD parental community should win more frequently. Signed $\Delta\mathrm{OD}_{A-B}$, the difference between the OD of parental community A and parental community B, should also be positively associated with PDI. We did not observe these predictions. The higher-OD parental community won only 37.1\% of Nutr$-$, 37.0\% of Base, and 13.2\% of Nutr$+$ Dominance events (Supplementary Fig.~45). Signed $\Delta\mathrm{OD}_{A-B}$ was not positively associated with PDI in any medium: it was negative in Nutr$-$, weak and not statistically significant in Base, and strongly negative in Nutr$+$ (Supplementary Fig.~14). OD-binned within-community pairwise selection correlation did not increase consistently across media (Supplementary Fig.~15). The negative association in Nutr$+$ (lower-OD communities winning more often) is consistent with the pH-mediated mechanism explored below: acidifying parental communities may yield lower endpoint biomass under high-nutrient conditions yet still win coalescence through environmental modification. Thus, parental-community biomass differences do not fully explain the nutrient-dependent increase in Dominance or correlated species fates.}

\subsubsection*{\rev{Initial richness and pool size}}
\rev{We also tested whether initial community richness could account for the observed community-level selection signatures in the experiments and simulations (Supplementary Fig.~13). In the synthetic-community experiments, we did not detect a statistically significant effect of inoculated isolate richness on the full three-category outcome distribution across 6, 12, and 24 inoculated isolates ($\chi^2 = 2.24$, $p = 0.69$). We also did not detect a pool-size effect within individual nutrient conditions (Dominance vs non-Dominance tests: Nutr$-$ $p = 0.39$, Base $p = 0.82$, Nutr$+$ $p = 0.27$). Initial inoculated richness affected assembly structure: realized parental ASV richness increased with pool size, whereas ASV richness per inoculated isolate decreased with pool size. These assembly-richness effects did not produce a corresponding pool-size dependence of Dominance frequency.}

\rev{In the gLV model, we performed a pool-size ablation at $\mu \in \{0.30, 0.60, 0.80\}$, varying the initial richness from 4 to 48 species per parental community while scaling the total pool as four non-overlapping parental communities ($N=16$--$192$ species). Dominance frequency was primarily set by $\mu$ and changed little with pool size, whereas the same-parental-community versus cross-parental-community pairwise selection correlation gap increased with pool size (Supplementary Fig.~13). This supports the interpretation that pool size affects assembly filtering and correlated species fates more than it affects Dominance-class frequency. Thus, pool-size variation does not fully explain the main transition in community-level selection; instead, it modulates the strength of correlated species fates within the interaction-strength regimes described in the main text.}

\subsubsection*{\rev{pH-feedback alternative model}}
\rev{We considered environmental-feedback models motivated by pH-mediated microbial interactions \citep{Ratzke2018,Ratzke2020}. A Ratzke--Gore-style pH-feedback model can generate asymmetric replacement over an intermediate feedback range, showing that environmental modification alone can contribute to Dominance-like outcomes (Supplementary Fig.~43). However, in our implementation this pH-only model did not reproduce the monotonic increase in Dominance with the interaction-strength parameter seen in the effective-interaction gLV framework and produced rare Restructuring outcomes ($\sim$1\%; 3/297 simulated events), in contrast to the higher Restructuring fraction in the effective-interaction gLV framework ($\sim$22\% across $\mu$ values). This result and the strict empirical pH-contrast analysis (Extended Data Fig.~7) support pH modification as a plausible contributor to winner identity in Nutr$+$, rather than as a complete explanation for the observed Nutr$+$ Dominance frequency or as a replacement for the outcome-level community-level selection interpretation.}

\subsubsection*{\rev{Strict pH contrast}}
\rev{In the strict acidic--alkaline subset, where one parental community had pH $<6.5$ and the other had pH $>7.5$, acidic parental communities dominated 29/32 Nutr$+$ events (90.6\%), compared with 23/41 Base events (56.1\%; Extended Data Fig.~7; Supplementary Fig.~46). A separate frequency comparison showed that strict acidic--alkaline pairings were enriched for Dominance relative to strict same-pH pairings in Nutr$+$ (29/32 versus 22/34 events; Fisher $p = 0.018$), but not in Base medium (24/41 versus 23/32 events; Fisher $p = 0.33$; Supplementary Fig.~46). This strict high-contrast analysis supports pH-mediated winner-direction bias in Nutr$+$, but it does not by itself explain the full Dominance pattern across all coalescence events because same-pH Nutr$+$ pairings often produced Dominance and intermediate-pH pairings were excluded.}

\subsubsection*{\rev{Dominant-species circularity}}
\rev{Dominant taxa could create circularity in the PDI analysis if the same species that drives a Dominance classification is also counted in the similarity calculation. We therefore recalculated PDI after removing either the most abundant species in the coalesced community or the dominant species from each parental community from all three composition vectors. The Nutr$+$ association weakened but remained statistically significant, whereas Base did not retain a statistically significant association (Supplementary Fig.~16). This supports the interpretation that dominant species contribute to the Nutr$+$ dominant-taxon/pH-mediated regime, but that the Nutr$+$ signal is not solely an artifact of including those same dominant species in the PDI calculation.}

\subsubsection*{\rev{Additive and geometric null effects}}
\rev{Abundance-skew and richness-dependent geometric effects can create apparent Dominance under simple null models. We address these effects in Supplementary Note~1, Extended Data Fig.~3, and Supplementary Fig.~38 using the coalescence-pair-specific simple additive null model and simulation parental-community composition-vector norm analyses. These analyses show that the simple additive null-model effect is present in principle, but does not account for most experimental coalescence outcomes and is insufficient to explain the simulated Dominance transition.}

\subsubsection*{\rev{Model-scope controls}}
\rev{The baseline gLV model does not represent facilitative interactions explicitly. The facilitative-tail, reciprocal pair-coupling, weak-mutualism, and coefficient mean/variance analyses are therefore retained in Supplementary Note~3 and Supplementary Figs.~39--42, where they define the scope of the effective-interaction framework. These analyses do not imply that facilitation is absent in all coalescence systems. They show that the main Dominance trend is not restricted to the narrowest iid competitive-matrix assumption.}
